% ══════════════════════════════════════════════════════════════════
%  FICHE D'EXERCICES — Chapitre 10 : La cinétique chimique
% ══════════════════════════════════════════════════════════════════
\section{Exercices type Baccalauréat}
\textit{La vitesse volumique de réaction se définit à partir de l'avancement ou des
concentrations. On rappelle $t_{1/2}$ : temps de demi-réaction.}

\rubrique{Vitesse de réaction}

\exo{}\diff{1} Définir la vitesse volumique de formation d'un produit et la vitesse
volumique de disparition d'un réactif. Comment les détermine-t-on graphiquement à
partir d'une courbe $[\text{X}] = f(t)$ ?

\exo{}\diff{1} Citer trois facteurs cinétiques. Pour chacun, indiquer le sens de son
influence sur la vitesse d'une réaction.

\exo{}\diff{2} On suit la réaction $\ce{X -> Y}$ par la concentration de $X$ :
\begin{center}\small
\begin{tabular}{l|cccccc}
$t$ (\si{\minute}) & 0 & 10 & 20 & 30 & 40 & 60 \\ \hline
$[X]$ (\si{\milli\mol\per\litre}) & 100 & 64 & 41 & 26 & 17 & 7 \\
\end{tabular}
\end{center}
\begin{enumerate}[label=\arabic*)]
\item Tracer l'allure de la courbe $[X]=f(t)$.
\item Déterminer la vitesse de disparition de $X$ à $t = 0$ et à $t = \SI{30}{\minute}$.
\item Comment évolue la vitesse au cours du temps ? Interpréter.
\item Déterminer le temps de demi-réaction $t_{1/2}$.
\end{enumerate}

\exo{}\diff{2} Lors d'une réaction, l'avancement $x$ passe de $0$ à
$x_{\max} = \SI{2.0}{\milli\mol}$. À $t = \SI{5}{\minute}$, l'avancement vaut
$x = \SI{1.0}{\milli\mol}$. Le volume du mélange est $V = \SI{100}{\milli\litre}$.
\begin{enumerate}[label=\arabic*)]
\item Que représente cet instant $t = \SI{5}{\minute}$ ?
\item Calculer la vitesse volumique moyenne de réaction entre $0$ et \SI{5}{\minute}.
\end{enumerate}

\rubrique{Facteurs cinétiques, catalyse}

\exo{}\diff{1} On réalise la même réaction à \SI{20}{\celsius} puis à \SI{50}{\celsius}.
Comparer qualitativement les durées de réaction et justifier.

\exo{}\diff{2} L'eau oxygénée $\ce{H2O2}$ se décompose selon
$\ce{2 H2O2 -> 2 H2O + O2}$. On ajoute des ions $\ce{Fe^3+}$.
\begin{enumerate}[label=\arabic*)]
\item Quel est le rôle des ions $\ce{Fe^3+}$ ? De quel type de catalyse s'agit-il ?
\item Comment pourrait-on suivre cette réaction expérimentalement ?
\end{enumerate}

\exo{}\diff{2} On étudie la réaction entre les ions iodure et l'eau oxygénée en
milieu acide : $\ce{H2O2 + 2 I- + 2 H+ -> I2 + 2 H2O}$. On suit la formation du
diiode par titrage.
\begin{enumerate}[label=\arabic*)]
\item Pourquoi peut-on suivre la réaction par dosage du diiode formé ?
\item Comment la vitesse est-elle modifiée si l'on augmente la concentration
initiale en ions iodure ? Justifier.
\end{enumerate}

\rubrique{Problème de synthèse — type Baccalauréat}

\sujetbac[d'après Baccalauréat C, Cameroun]
On étudie la cinétique de la réaction $\ce{S2O8^2- + 2 I- -> 2 SO4^2- + I2}$. On
mesure la concentration en diiode $[\ce{I2}]$ formé au cours du temps :
\begin{center}\small
\begin{tabular}{l|ccccccc}
$t$ (\si{\minute}) & 0 & 2 & 4 & 6 & 8 & 12 & 20 \\ \hline
$[\ce{I2}]$ (\si{\milli\mol\per\litre}) & 0 & 3{,}5 & 6{,}0 & 7{,}6 & 8{,}6 & 9{,}5 & 10{,}0 \\
\end{tabular}
\end{center}
\begin{enumerate}[label=\arabic*)]
\item Tracer la courbe $[\ce{I2}] = f(t)$.
\item Définir et déterminer la vitesse de formation du diiode à $t = 0$ et à
$t = \SI{8}{\minute}$.
\item Comment évolue cette vitesse ? Quel facteur cinétique l'explique ?
\item Déterminer la concentration finale en diiode et le temps de demi-réaction.
\item On recommence l'expérience à une température plus élevée. Tracer l'allure de la
nouvelle courbe sur le même graphe et justifier.
\end{enumerate}
